Si può dimostrare che le operazioni su $\mathbb{R}^n$ definite nella sezione precedente godono
delle seguenti proprietà.

\begin{property}
    \textbf{Commutatività della somma su $\mathbb{R}^n$}
    \[ \forall\overline{x},\overline{y}\in\mathbb{R}^n,
    \overline{x}+\overline{y}=\overline{x}+\overline{y}. \]
\end{property}
\begin{property}
    \textbf{Associatività della somma su $\mathbb{R}^n$}
    \[ \forall\overline{x},\overline{y},\overline{z}\in\mathbb{R}^n,
    \left(\overline{x}+\overline{y}\right)+\overline{z}=
    \overline{x}+\left(\overline{y}+\overline{z}\right). \]
\end{property}
\begin{property}
    \textbf{Esistenza dell'elemento neutro per la somma su $\mathbb{R}^n$}
    \[ \forall \overline{x}\in\mathbb{R}^n,  \exists\overline{e}=\left(0, 0,
    \ldots, 0\right)=\overline{0}\in\mathbb{R}^n \mid \overline{x}+\overline{e}
    =\overline{e}+\overline{x}=\overline{x}.\]
\end{property}
\begin{property}
    \textbf{Esistenza dell'elemento inverso per la somma}sss
    \[ \forall \overline{x}\in\mathbb{X},  \exists\overline{y}=\left(-x_1, -x_2,
    \ldots, -x_n\right)=-\overline{x}\in\mathbb{X} \mid \overline{x}+\overline{y}
    =\overline{y}+\overline{x}=\overline{e}.\]
\end{property}


\begin{property}
    \textbf{Commutatività del prodotto per scalare}
    \[ \forall\overline{x}\in\mathbb{R}^n, \,\alpha\in\mathbb{R},\;
    \alpha \overline{x} = \overline{x} \alpha \]
\end{property}
\begin{property}
    \textbf{Associatività del prodotto per scalare}
    \[ \forall\overline{x}\in\mathbb{R}^n,\, \alpha,\beta\in\mathbb{R},\;
    \alpha\left(\beta \overline{x}\right) = \left(\alpha\beta\right)\overline{x}. \]
\end{property}
\begin{property}
    \textbf{Distributività del prodotto per scalare sulla somma su $\mathbb{R}^n$}
    \[ \forall\overline{x},\overline{y}\in\mathbb{R}^n,\, \alpha\in\mathbb{R},\;
    \alpha\left(\overline{x}+\overline{y}\right) = \alpha\overline{x}+\alpha\overline{y}. \]
\end{property}


\begin{property}\label{p:pro-sca-1}
    \textbf{Proprietà 1 del prodotto scalare euclideo}
    \[ \forall\overline{x},\overline{y}\in\mathbb{R}^n,\;
    \langle\overline{x},\overline{y}\rangle \in\mathbb{R}. \]
\end{property}

\begin{property}\label{p:pro-sca-2}
    \textbf{Proprietà 2 del prodotto scalare euclideo}
    \[ \forall\overline{x}\in\mathbb{R}^n,\;
    \langle\overline{x},\overline{x}\rangle \geq 0. \]
    Inoltre
    \[ \langle\overline{x},\overline{x}\rangle = 0 \Leftrightarrow \overline{x}=\overline{0}. \]
\end{property}

\begin{property}\label{p:pro-sca-3}
    \textbf{Proprietà 3 del prodotto scalare euclideo}
    \[ \forall\overline{x},\overline{y}\in\mathbb{R}^n,\;
    \langle\overline{x},\overline{y}\rangle=\langle\overline{y},\overline{x}\rangle. \]
\end{property}

\begin{property}\label{p:pro-sca-4}
    \textbf{Proprietà 4 del prodotto scalare euclideo (additività)}
    \[ \forall\overline{x},\overline{y},\overline{z}\in\mathbb{R}^n,\;
    \langle\overline{x}+\overline{y},\overline{z}\rangle=\langle\overline{x},\overline{z}\rangle + \langle\overline{y},\overline{z}\rangle. \]
\end{property}

\begin{property}\label{p:pro-sca-5}
    \textbf{Proprietà 5 del prodotto scalare euclideo (omogeneità)}
    \[ \forall\overline{x},\overline{y}\in\mathbb{R}^n,\, \alpha\in\mathbb{R},\;
    \alpha\langle\overline{x},\overline{y}\rangle = \langle\alpha\overline{x},\overline{y}\rangle = \langle\overline{x},\alpha\overline{y}\rangle. \]
\end{property}